% Appendix D -- Procedurer

\chapter{Procedurer}
\label{AppendixD}

Dette bilag samler de praktiske procedurer, der bruges til at
opstarte og fejlsøge robotten i den konfiguration, som rapporten
omfatter. Procedurerne er en opdatering af det eksisterende
materiale i workspace-repositoriet.

% ---------------------------------------------------------------------
\section{Opstartsprocedure}
\label{sec:appC-opstart}

Den følgende procedure tager udgangspunkt i den afgrænsede
konfiguration, hvor kun venstre arm er aktiv. Trinene udføres på
Jetson Orin Nano-hosten.

\begin{enumerate}
    \item \textbf{Forberedelse.} Tjek, at de tre ESP32-kontrollere
        sidder i de korrekte USB-porte (rød i port~2.2, gul i
        port~2.3, hvid i port~2.1, jf.~afsnit~\ref{sec:appC-portbinding}),
        og at strømforsyningen til armen er tilsluttet og tændt.
    \item \textbf{Start manager-noden.} Kør
        \enquote{ros2 launch wattson\_servos servos.launch.py}.
        Manager-noden indlæser
        \enquote{servo\_arm\_left\_params.json}, åbner forbindelsen
        til ESP32'en for venstre arm via dens \enquote{by-path}-sti,
        og opretter ROS-topics og services for de syv led.
    \item \textbf{Start slider-GUI'en.} Kør
        \enquote{ros2 launch wattson\_servos slider\_control.launch.py}.
        GUI'en åbner med en skyder pr.~led, så hvert led kan
        bevæges uafhængigt.
    \item \textbf{Verifikation.} Bevæg én skyder ad gangen, og
        bekræft visuelt, at det rigtige led reagerer i den rigtige
        retning og inden for sine software-grænser
        (jf.~Tabel~\ref{tab:kalibrering-tab} i
        Bilag~\ref{AppendixB}).
\end{enumerate}

Forsyningen skal være slukket, inden kabler eller stik flyttes,
og inden manager-noden genstartes mod en ny konfiguration.

% ---------------------------------------------------------------------
\clearpage
\section{Portbinding -- procedure}
\label{sec:appC-portbinding}

Dette afsnit beskriver, hvordan den faste sti i
\enquote{/dev/serial/by-path/} for hver ESP32-kontroller findes,
og hvordan de tre stier hænger sammen med kabelfarverne på den
fysiske montage.

\paragraph{Find stien for en kontroller.}
\begin{enumerate}
    \item Tilslut én ESP32 ad gangen og kør
        \enquote{ls -l /dev/serial/by-path/}. Stien til den
        netop tilsluttede kontroller dukker op som et symlink, der
        peger ind på den tilhørende \enquote{/dev/ttyUSBx}.
    \item Noter sti og fysisk port (rød/gul/hvid) for hver
        kontroller.
    \item Indsæt stierne i manager-nodens konfiguration som
        beskrevet i Bilag~\ref{AppendixC},
        afsnit~\ref{sec:appB-portbinding}.
\end{enumerate}

\paragraph{Mapping mellem port, kabelfarve og delsystem.}
Den verificerede mapping for den nuværende fysiske montage er
gengivet i Tabel~\ref{tab:portmapping}.

\begin{table}[H]
\centering
\caption{Port-til-kabel-mapping for de tre ESP32-kontrollere
på Jetson Orin Nano.}
\label{tab:portmapping}
\begin{tabular}{@{}lll@{}}
\toprule
\textbf{Kabelfarve} & \textbf{USB-port} & \textbf{Delsystem} \\
\midrule
Rød  & 2.2 & Venstre arm \\
Gul  & 2.3 & Højre arm + hoved \\
Hvid & 2.1 & Begge hænder \\
\bottomrule
\end{tabular}
\end{table}
